\begin{corollary}[非原子冻结的正熵凝聚]\label{cor:conclusion-realinput40-nonatomic-freezing-positive-entropy}
\leanverified{paper\_conclusion\_realinput40\_nonatomic\_freezing\_positive\_entropy}
沿用小节 \ref{subsec:conclusion-realinput40-analytic-window-logshoulder-ground-quartic-unification} 的记号。则真实输入 \(40\) 状态核在饱和端点的冻结并不塌缩为单一轨道或原子态，而是凝聚到一个四状态、正熵、有限型的地基 SFT。更精确地，
\[
\alpha_{\max}=\lim_{\theta\to+\infty}P'(\theta)=\frac12,
\]
且在缩放 \(z=\sqrt{w/u}\) 下，
\[
\Delta\!\left(\sqrt{w/u},u\right)\longrightarrow (1-w)\det(I-wB_\infty),
\]
同时
\[
h_{\mathrm{ground}}=\frac12\log\rho(B_\infty)=\log c>0.
\]
因此原子因子 \((1-w)\) 只留下可剥离的边界残余，而主导冻结尺度的是四状态地基层 \(\det(I-wB_\infty)\)。
\end{corollary}
\begin{proof}
由推论 \ref{cor:real-input-40-alpha-max}，
\[
\alpha_{\max}=\frac12.
\]
再由定理 \ref{thm:conclusion-realinput40-zero-temp-noncrystallization}，
\[
\lim_{u\to+\infty}\Delta\!\left(\sqrt{w/u},u\right)
=(1-w)\det(I-wB_\infty),
\]
且
\[
\lim_{n\to\infty}\frac1n\log A_{n,k_{\max}(n)}=\log c>0.
\]
这正说明零温极限并不结晶到有限条周期轨道，而是保留一个正熵四状态地基 SFT。最后，由定理 \ref{thm:conclusion-realinput40-zero-temp-quartic-constant-unification}，
\[
b=\rho(B_\infty)^{-1}\in(0,1),
\]
故最内层主导零点来自 \(\det(I-wB_\infty)\) 的正根 \(w=b\)，而不是边界因子 \(w=1\)。证毕。
\end{proof}

\begin{theorem}[镜像 Perron 驱动的平方根慢化]\label{thm:conclusion-realinput40-mirror-perron-sqrt-slowdown}
沿用正文饱和端点的 Puiseux 记号。除显式谱支
\[
\lambda_\pm(u)=\pm\sqrt u
\]
外，还存在一条负实镜像 Perron 分支 \(\lambda_{\mathrm{mir}}(u)\)，满足
\[
\lambda_1(u)=c\sqrt u+d+\frac{e}{\sqrt u}+O(u^{-1}),
\qquad
\lambda_{\mathrm{mir}}(u)=-c\sqrt u+d-\frac{e}{\sqrt u}+O(u^{-1}),
\]
从而
\[
\Lambda_2(u)=|\lambda_{\mathrm{mir}}(u)|\sim c\sqrt u,
\]
并且
\[
\boxed{\;
1-\frac{\Lambda_2(u)}{\lambda_1(u)}
=
\frac{2d}{c}u^{-1/2}+O(u^{-1}).\;}
\]
于是由谱隙比控制的弛豫时间满足
\[
\boxed{\;\tau_{\mathrm{rel}}(u)\sim \frac{c}{2d}\sqrt u.\;}
\]
更强地，对任意固定 \(\beta<1\)，存在 \(u_0(\beta)\) 使得当 \(u\ge u_0(\beta)\) 时，
\[
\Lambda_2(u)>\lambda_1(u)^\beta.
\]
因此，在饱和端点邻域不存在统一的幂指数型次谱节约；真正主导慢化的不是显式 \(\pm\sqrt u\) 分支，而是与主 Perron 根发生近碰撞的镜像 Perron 分支。
\end{theorem}
\begin{proof}
正文已给出主 Perron 根与镜像 Perron 分支的 Puiseux 展开。由于 \(\lambda_{\mathrm{mir}}(u)<0\) 对充分大的 \(u\) 成立，取绝对值即得
\[
\Lambda_2(u)=|\lambda_{\mathrm{mir}}(u)|=c\sqrt u-d+\frac{e}{\sqrt u}+O(u^{-1}),
\]
从而
\[
\frac{\Lambda_2(u)}{\lambda_1(u)}
=
\frac{c\sqrt u-d+e\,u^{-1/2}+O(u^{-1})}{c\sqrt u+d+e\,u^{-1/2}+O(u^{-1})}
=
1-\frac{2d}{c}u^{-1/2}+O(u^{-1}).
\]
整理即得
\[
1-\frac{\Lambda_2(u)}{\lambda_1(u)}
=
\frac{2d}{c}u^{-1/2}+O(u^{-1}).
\]
若以 \((1-\Lambda_2/\lambda_1)^{-1}\) 作为弛豫时间标度，则
\[
\tau_{\mathrm{rel}}(u)\sim \frac{c}{2d}\sqrt u.
\]

最后，因
\[
\Lambda_2(u)\sim c\,u^{1/2},
\qquad
\lambda_1(u)^\beta\sim c^\beta u^{\beta/2},
\]
而 \((1-\beta)/2>0\)，故
\[
\frac{\Lambda_2(u)}{\lambda_1(u)^\beta}\sim c^{1-\beta}u^{(1-\beta)/2}\to+\infty.
\]
因此对充分大的 \(u\) 必有 \(\Lambda_2(u)>\lambda_1(u)^\beta\)。证毕。
\end{proof}
